% SPDX-License-Identifier: GPL-3.0-or-later
% Copyright (C) 2026 Alessandro Greco
\section*{Глава \ref{cap:01-anonimato}: анонимность, наблюдение и метаданные}

\begin{soluzione}{es:01-metadati}
При HTTPS интернет-провайдер видит IP-адрес веб-сайта и потенциально его доменное имя — из DNS-разрешения и часто из поля TLS SNI. Он также видит временные характеристики, длительность соединения и объём трафика. При использовании VPN провайдер видит только соединение с VPN-сервером; сам VPN-провайдер, однако, может видеть всё то, что раньше видел интернет-провайдер. При использовании Tor провайдер видит соединение с ретранслятором Tor или мостом, а также его временные характеристики и объём. Ни одна отдельная точка внутри сети Tor не видит одновременно пользователя и назначение.
\end{soluzione}

\begin{soluzione}{es:01-termini}
Учётная запись под вымышленным именем — пример псевдонимности. Неподписанное сообщение на доске объявлений — пример анонимности, если ничто другое не идентифицирует автора. Письмо, подписанное инициалами, — слабая псевдонимность и нередко позволяет установить личность. Анонимное пожертвование наличными — пример анонимности.
\end{soluzione}

\begin{soluzione}{es:01-insieme}
Множество анонимности сокращается от пятидесяти пользователей Tor до тех, кто говорит по-итальянски и одновременно бодрствует в 3 часа ночи; потенциально это всего несколько человек. Tor скрывает IP-адрес, а не поведение. Язык, время, стиль письма и личные сведения способны идентифицировать пользователя даже тогда, когда сетевой трафик этого не позволяет.
\end{soluzione}

\begin{soluzione}{es:01-vpn}
VPN-провайдер видит и пользователя, и его назначения. Всё доверие сосредоточено у одного оператора, который может вести журналы, быть юридически обязанным передать данные или быть скомпрометирован. Tor распределяет доверие между тремя независимо управляемыми ретрансляторами, ни один из которых не видит оба конца. Для связывания пользователя и назначения пришлось бы скомпрометировать или одновременно наблюдать первый и последний переход.
\end{soluzione}

\section*{Глава \ref{cap:02-radici}: истоки}

\begin{soluzione}{es:02-mix}
При случайном выборе вероятность составляет 1 из 100. Если наблюдатель сам отправил 99 из 100 сообщений в партии, он способен распознать собственные исходящие сообщения и определить единственное оставшееся сообщение как сообщение жертвы. Анонимность исчезает. Это \emph{атака $n-1$}; реальные сети смешивания противодействуют ей более сложным формированием партий, фиктивными сообщениями и контролем поступающих сообщений.
\end{soluzione}

\begin{soluzione}{es:02-penet}
Используйте \emph{анонимный обратный адрес} в смысле Чаума. Пользователь подготавливает зашифрованный многослойный \enquote{обратную луковицу} для цепочки узлов смешивания. Внутренний слой содержит адрес пользователя и ключ для шифрования ответа. Отвечающий присоединяет сообщение к луковице; каждый узел смешивания снимает один слой, и ни один отдельный узел смешивания не узнаёт конечное назначение. Центральная таблица маршрутизации не требуется.
\end{soluzione}

\begin{soluzione}{es:02-crowds}
На каждом шаге запрос пересылается другому участнику группы с вероятностью $3/4$ и отправляется веб-сайту с вероятностью $1/4$. Число пройденных участников имеет геометрическое распределение со средним $1/(1/4) = 4$, если учитывать первый шаг пересылки. Участник не может знать, является ли переславший запрос узел его источником, поскольку промежуточная пересылка подчиняется тому же правилу.
\end{soluzione}

\section*{Глава \ref{cap:03-nascita-tor}: от onion routing NRL к Tor Project}

\begin{soluzione}{es:03-cinque}
При трёх переходах вредоносный центральный узел сразу знал две организационные сети на концах, поскольку входной и выходной узлы находились внутри сетей этих организаций. При четырёх переходах тот, кто скомпрометировал узел рядом с одним концом, знал, какой дополнительный узел атаковать для восстановления цепи; при пяти переходах требовалась ещё одна компрометация. Tor использует три перехода, поскольку его ретрансляторы не находятся внутри сетей пользователей и потому что каждый дополнительный переход увеличивает задержку, не уменьшая существенно главную угрозу корреляции входом и выходом. Вместо этого против этой угрозы помогают входные узлы (guard).
\end{soluzione}

\begin{soluzione}{es:03-cipolla}
Единая луковица, зашифрованный долгосрочными открытыми ключами узлов, не обеспечивает прямую секретность. Тот, кто позднее получит эти ключи, сможет расшифровать ранее записанные луковицы. При телескопическом построении на каждом переходе выполняется эфемерный обмен ключами, а временные секреты после использования удаляются. Клиент также строит цепь по одному переходу и может сразу определить не отвечающий ретранслятор и заменить его. При единой луковице отказ становится заметен только после попытки пройти весь маршрут.
\end{soluzione}

\begin{soluzione}{es:03-fiducia}
Это упражнение с открытым аргументированным ответом. К проверяемым элементам относятся опубликованные исходный код и спецификации, воспроизводимые сборки, независимый анализ в научной литературе и сеть, управляемая независимыми операторами. Доверия требуют, например, соответствие распространяемых сборок опубликованному исходному коду, хотя это частично проверяемо; отсутствие сговора центров каталогов; отсутствие неизвестных уязвимостей. Безопасность следует оценивать по проверяемости и распределённому доверию, а не выводить из источников финансирования проекта.
\end{soluzione}

\section*{Глава \ref{cap:04-modello-minaccia}: модель угроз}

\begin{soluzione}{es:04-casi}
(a) Да: оператор сети видит только соединение с ретранслятором Tor. (b) Да: веб-сайт видит адрес выходного ретранслятора. (c) Нет в отношении идентичности: пользователь сам идентифицирует себя, входя в учётную запись. Tor скрывает только место, откуда выполняется соединение. (d) Вероятно, нет: противник, наблюдающий оба конца — пользователя и газету, — может попытаться выполнить временную корреляцию, то есть именно тот сценарий, который исключён из модели угроз Tor.
\end{soluzione}

\begin{soluzione}{es:04-usabilita}
Сложная система привлекает лишь небольшое число сильно мотивированных пользователей, часто именно тех, кому защита нужна больше всего. Уже само её использование выделяет их и оставляет небольшое множество анонимности. Простая в использовании система, которой миллионы людей пользуются по обычным причинам, помогает скрываться тем, чьи потребности более чувствительны.
\end{soluzione}

\begin{soluzione}{es:04-modello}
Это упражнение с открытым ответом. Возможные противники включают государство или работодателя источника, наблюдающих локальную сеть, провайдера электронной почты или сообщений и организации, контролирующие сети в одной из двух стран. Tor не позволяет локальной сети узнать, с кем общается журналист, а сервису — увидеть IP-адрес журналиста. Он не защищает содержимое сообщений без сквозного шифрования, метаданные, видимые сервису, например учётные записи и временные метки, скомпрометированные устройства или корреляцию трафика противником, наблюдающим оба конца. Onion-сервис вроде SecureDrop уменьшает несколько из этих рисков.
\end{soluzione}

\section*{Глава \ref{cap:05-crittografia}: необходимая криптография}

\begin{soluzione}{es:05-xor}
$C_1 \oplus C_2 = P_1 \oplus P_2$, поскольку повторно использованный поток ключа сокращается. Зная $P_1$, можно восстановить $P_2 = C_1 \oplus C_2 \oplus P_1$.
\end{soluzione}

\begin{soluzione}{es:05-shake}
16-байтовый результат совпадает с первыми 16 байтами 32-байтового результата: SHAKE генерирует один поток, который можно считывать до любой требуемой длины. Поэтому KDF может запросить длинный поток и разделить его на части. Однако запрос различных длин для разных целей \emph{не} обеспечивает разделения доменов; для этого всё равно нужны разные доменные метки.
\end{soluzione}

\begin{soluzione}{es:05-dh}
Оба общих секрета совпадают. Полностью нулевой открытый ключ представляет точку малого порядка и дал бы полностью нулевой общий секрет независимо от закрытого ключа, то есть секрет, известный всем. Библиотека \texttt{cryptography} отказывается его возвращать. ntor явно требует эту проверку (\cref{sec:ntor}); без неё атакующий мог бы навязать предсказуемый общий секрет.
\end{soluzione}

\begin{soluzione}{es:05-etichette}
Без разделения доменов значение, вычисленное в одном контексте, может быть принято в другом. Например, MAC, созданный для аутентификации сообщения рукопожатия, можно попытаться представить как ключ или доказательство в другой фазе либо в другом протоколе, использующем ту же функцию и ключи. Доменные метки по конструкции делают входы разных назначений непересекающимися.
\end{soluzione}

\section*{Глава \ref{cap:06-architettura}: архитектура сети}

\begin{soluzione}{es:06-vista}
Входной узел (guard): IP-адрес пользователя, временные характеристики и объём трафика, но не назначение. Промежуточный ретранслятор: ни пользователя, ни назначения. Выходной ретранслятор: назначение и любое содержимое, не зашифрованное приложением, но не пользователя. HSDir: никакой информации о пользователе, посещающем обычный веб-сайт. Центр каталогов: обычно никакой информации о конкретном посещении сайта. Она может видеть IP-адрес клиента, загружающего документы каталогов \emph{напрямую}, особенно при первоначальном подключении. Если клиент загружает их через резервный сервер каталогов или кэш, запрос видит этот сервер, а не центр каталогов.
\end{soluzione}

\begin{soluzione}{es:06-pr}
\texttt{Link=3-5} объявляет версии протокола связи вплоть до версии 5 (\cref{cap:07-canali}). \texttt{Relay=2-6} включает ntor, ntor-v3 (\texttt{Relay=4}) и CGO (\texttt{Relay=6}). \texttt{FlowCtrl=1-2} означает аутентифицированные SENDME и управление перегрузкой. Ретранслятор может выступать выходным узлом с CGO при условии, что также имеет флаг \texttt{Exit} и согласует управление перегрузкой вместе с CGO (\cref{sec:ntor3}).
\end{soluzione}

\begin{soluzione}{es:06-livelli}
100\,KB прикладных данных требуют примерно $102400/498 \approx 206$ ячеек DATA, то есть около $206 \times 514 \approx 106$\,KB на канале. Нужно добавить несколько ячеек для запроса, открытия потока и сообщений SENDME — примерно семь, по одному на каждые 31 ячейку данных, — а также служебных данных записей TLS. В сумме это примерно на 5--7\% больше полезной нагрузки, без учёта начальных рукопожатий.
\end{soluzione}
